\subsection{Choice, Zorn's Lemma, and Well-Ordering}

\begin{principle} \label{princ:axiom-of-choice-analysis} \hlt{Choice Principle}\\
For nonempty sets $\{X_i\}_{i\in I}$, the axiom of choice asserts $\exists c:I\to\bigcup\limits_{i\in I}X_i$ such that $c(i)\in X_i\ \forall i\in I$.\\
Use direct construction for finite or explicit choices; use choice only for arbitrary simultaneous choices.
\end{principle}

\begin{principle} \label{princ:well-ordering-natural-numbers} \hlt{Well-Ordering of $\N$}\\
Every nonempty subset of $\N$ has a least element, justifying first unused indices, first repeated values, and increasing listings of infinite subsets of $\N$.
\end{principle}

\begin{principle} \label{princ:zorns-lemma} \hlt{Zorn's Lemma}\\
If a nonempty partially ordered set $X$ has the property that every chain in $X$ has an upper bound in $X$, then $X$ contains a maximal element. This is equivalent in strength to the axiom of choice, Principle \ref{princ:axiom-of-choice-analysis}.
\end{principle}

\hlt{Placeholder}: Add the well-ordering theorem, maximality methods, and examples of where choice is avoidable versus structurally necessary.
